\documentclass[10pt,a4paper]{amsart}
\usepackage[margin=19mm]{geometry}
\usepackage{amsmath,amssymb,mathtools}
\usepackage[scr=rsfs,cal=euler]{mathalfa}
\usepackage{xfrac}
\usepackage[unicode,hidelinks,bookmarks=false]{hyperref}
\newcommand{\ie}{i.e.\ }
\newcommand{\cf}{cf.\ }
\newcommand{\resp}{respectively\ }
\newcommand{\Subsection}{\subsection}
\providecommand{\nobreakdash}{\nobreak\hbox{-}}
\setlength{\emergencystretch}{2em}
\multlinegap0pt
\usepackage{bm}
\usepackage{bbm}
\usepackage{dsfont}
\usepackage{xspace}
\usepackage{cancel}
\usepackage{mathtools}
\usepackage{amsthm}
\makeatletter
\@ifundefined{theorem}{\newtheorem{theorem}{Theorem}}{}
\@ifundefined{definition}{\newtheorem{definition}[theorem]{Definition}}{}
\@ifundefined{lemma}{\newtheorem{lemma}[theorem]{Lemma}}{}
\makeatother
\title[An Extrapolation Theorem in the setting of Hausdorff capacit]{An Extrapolation Theorem in the setting of Hausdorff capacities}
\author{Aniruddha Deshmukh}
\date{Source: arXiv:2510.04105v1 (CC BY 4.0)}
\begin{document}
\maketitle

\noindent\emph{Source and licence.} An Extrapolation Theorem in the setting of Hausdorff capacities. Version arXiv:2510.04105v1 (CC BY 4.0), distributed under the Creative Commons Attribution 4.0 International licence, as recorded in the arXiv licence metadata for this exact version and shown on the abstract page https://arxiv.org/abs/2510.04105v1. Full rights notice: licenses/arxiv-2510.04105-CC-BY-4.0.txt.\par\smallskip

\noindent\textit{Editorial scope.} Counted unit: the Theorem whose source label is \texttt{ConstructionA1Weights} in the article (source0.tex lines 319--322, source label \texttt{ConstructionA1Weights}), delivered together with the article's own complete original proof of it (source0.tex lines 323--356, one \texttt{proof} environment of 34 lines). No mathematics is written, repaired or re-derived: statement and proof are the article's own text line for line. Required context retained uncounted: CubicHC (lines 105--112); PreparatoryLemmaConstructionA1 (lines 294--306). Every definition, display and label of those lines is retained unchanged. Omitted from this package and not redistributed: 1--104, 113--293, 307--318, 357--436. Nothing inside a retained excerpt is omitted or altered: every retained body is delivered exactly as the article's lines are written, so no omission receipt of any kind accompanies this package. Citations: the retained excerpts carry no citation command at all, so no bibliography entry of the article is transcribed and no reference is redistributed. Reference adaptation: none was needed and none was made; every reference inside the delivered unit resolves inside it, so no unresolved reference or citation is produced. Printed slips kept literal: no printed word, symbol, hypothesis or number of the counted statement or of its proof was altered. Layout and class: the article's own class is \texttt{amsart} with packages including amsthm, amsmath, amssymb, enumerate, enumitem, mathrsfs, float, graphicx, mathtools, hyperref, caption, subcaption, faktor, xcolor; the wrapper is the lane's pinned \texttt{amsart} editorial wrapper at 10pt on A4, which restates the theorem-like environments the retained text needs and adds the article's own macro definitions that the retained text needs (none), with extra wrapper packages (bm, bbm, dsfont, xspace, cancel, mathtools). The delivered numbering is the unit's own. Assets: no retained span contains a figure environment, a graphics-inclusion command, a PDF-page inclusion, a table or a TikZ picture; no image file is copied into this package, the package declares no asset dependency and no image file is redistributed with it. Licence: version arXiv:2510.04105v1 of this article is distributed under the Creative Commons Attribution 4.0 International licence, as recorded in the arXiv licence metadata for this exact version and shown on the abstract page https://arxiv.org/abs/2510.04105v1. A full rights notice is delivered at licenses/arxiv-2510.04105-CC-BY-4.0.txt. Characters: every retained excerpt is pure ASCII, so the delivered engine, XeLaTeX with its default Latin Modern text font, reports no missing character.
		\begin{definition}[Cubic Hasudorff Content]
			\label{CubicHC}
			The cubic Hausdorff content is a set function $H_\infty^{\beta,c}: \mathcal{P}(\mathbb{R}^n) \rightarrow [0,\infty],$ defined as
			\begin{equation}
				H_\infty^{\beta,c} (E) := \inf\left \{\sum_{j=1}^{\infty}\ell (Q_j)^\beta : E \subseteq \bigcup_{i=1}^\infty Q_j \right\},
			\end{equation}
			where, $Q_j$ is a cube in $\mathbb{R}^n$ with sides parallel to the coordinate axes and side length $ \ell(Q_j)$. 
		\end{definition}

		\begin{lemma}
			\label{PreparatoryLemmaConstructionA1}
			Let $T$ be an operator acting on functions on $\mathbb{R}^n$ such that there is a constant $K > 0$ such that for every $f \in L^1 \left( \mathbb{R}^n, H_{\infty}^{\beta} \right)$, we have,
			\begin{equation}
				\label{Weak11T}
				H_{\infty}^{\beta} \left( \left\lbrace x \in \mathbb{R}^n | \left| Tf \left( x \right) \right| > t \right\rbrace \right) \leq \frac{K}{t} \| f \|_{L^1 \left( \mathbb{R}^n, H^{\beta}_{\infty} \right)}.
			\end{equation}
			Then, For every $0 < \gamma < 1$, there is a constant $K' > 0$ depending only on $\gamma$ such that for any set $E \subseteq \mathbb{R}^n$ with $H^{\beta}_{\infty} \left( E \right) < + \infty$, we have,
			\begin{equation}
				\label{PreparatoryEquation}
				\int\limits_{E} \left| Tf \right|^{\gamma} \mathrm{d}H^{\beta}_{\infty} \leq K' \left( H^{\beta}_{\infty} \left( E \right) \right)^{1 - \gamma} \| f \|_{L^1 \left( \mathbb{R}^n, H^{\beta}_{\infty} \right)}^{\gamma}.
			\end{equation}
		\end{lemma}

		\begin{theorem}
			\label{ConstructionA1Weights}
			Let $f \in L^1_{\text{loc}} \left( H^{\beta}_{\infty} \right)$ such that $M_{\beta}f \left( x \right) < + \infty$ for $H^{\beta}_{\infty}$-quasi-every $x \in \mathbb{R}^n$. Then, for every $\delta \in \left[ 0, 1 \right)$, the function $\left( M_{\beta}f \right)^{\delta}$ is an $A_{1, \beta}$ weight.
		\end{theorem}

		\begin{proof}
			Given $x \in \mathbb{R}^n$, we fix a cube $Q \subseteq \mathbb{R}^n$ with $x \in Q$. To show that $\left( M_{\beta}f \right)^{\delta}$ satisfies the $A_{1, \beta}$ condition, it is enough to show that with $Q$ fixed arbitrarily, we have,
			\begin{equation}
				\label{GoalEquation}
				\frac{1}{H^{\beta}_{\infty} \left( Q \right)} \int\limits_{Q} \left( M_{\beta}f \right)^{\delta} \ \mathrm{d}H^{\beta}_{\infty} \leq K \left( M_{\beta}f \right)^{\delta} \left( x \right),
			\end{equation}
			for some constant $K > 0$. In order to do so, we decompose $f$ as
			$$f = f \cdot \chi_{2Q} + f \cdot \chi_{\mathbb{R}^n \setminus \left( 2Q \right)},$$
			where, $2Q$ is the cube with the same center as that of $Q$ but twice its length. For the simplicity of notations, let us write $f_1 = f \cdot \chi_{2Q}$ and $f_2 = f \cdot \chi_{\mathbb{R}^n \setminus \left( 2Q \right)}$. Since the Choquet integral for the dyadic Hausdorff content $H^{\beta, Q_0}_{\infty}$ is sublinear and the dyadic Hausdorff content is equivalent to $H^{\beta}_{\infty}$, we have,
			$$M_{\beta}f \left( y \right) \leq K \left( M_{\beta}f_1 \left( y \right) + M_{\beta}f_2 \left( y \right) \right),$$
			for some constant $K > 0$ that depends only on the dimensions $n$ and $\beta$. Further, since $0 \leq \delta < 1$, we see that
			$$\left( M_{\beta}f \right)^{\delta} \left( y \right) \leq K^{\delta} \left( \left( M_{\beta}f_1 \right)^{\delta} \left( y \right) + \left( M_{\beta}f_2 \right)^{\delta} \left( y \right) \right),$$
			for every $y \in \mathbb{R}^n$. From this point onward, we use the alphabet $K$ to denote a positive constant that depends only on the parameters of the statement, and whose value might be different at every occurrence.
			
			Since $M_{\beta}$ is satisfies the weak type $(1, 1)$ inequality with respect to $H^{\beta}_{\infty}$, we have from Lemma \ref{PreparatoryLemmaConstructionA1},
			$$\frac{1}{H^{\beta}_{\infty} \left( Q \right)} \int\limits_{Q} \left( Mf_1 \right)^{\delta} \mathrm{d}H^{\beta}_{\infty} \leq K \left( \frac{1}{H^{\beta}_{\infty} \left( Q \right)} \int\limits_{2Q} \left| f_1 \right| \mathrm{d}H^{\beta}_{\infty} \right)^{\delta}.$$
			Now, we interpret (upto a constant) $H^{\beta}_{\infty}$ as the cubic Hausdorff content (see Definition \ref{CubicHC}) so that $H^{\beta}_{\infty} \left( Q \right) = \left( \ell \left( Q \right) \right)^{\beta} = \frac{H^{\beta}_{\infty} \left( 2Q \right)}{2^{\beta}}$. That is, we have,
			\begin{equation}
				\label{F1Equation}
				\frac{1}{H^{\beta}_{\infty} \left( Q \right)} \int\limits_{Q} \left( Mf_1 \right)^{\delta} \mathrm{d}H^{\beta}_{\infty} \leq K \left( \frac{1}{H^{\beta}_{\infty} \left( 2Q \right)} \int\limits_{2Q} \left| f \right| \mathrm{d}H^{\beta}_{\infty} \right)^{\delta} \leq K \left( M_{\beta}f \right)^{\delta} \left( x \right).
			\end{equation}
			Next, we take a cube $Q'$ containing a point $y \in Q$ such that $\int\limits_{Q'} \left| f_2 \right| \mathrm{d}H^{\beta}_{\infty} > 0$. Since $f_2$ is supported outside of $2Q$, we must have that the length of $Q'$ is at least half that of $Q$, for if not $Q' \subseteq 2Q$ and the integral would be zero. This also gives that $2Q \subseteq K \cdot Q'$, for a constant $K > 0$ that depends only on the dimension $n$ (for instance, one may take $K = 4 \sqrt{n}$). This gives that $K \cdot Q$ contains $x$.
			
			Thus, for any $y \in Q$ and a cube $Q'$ containing $y$ with $\int\limits_{Q'} \left| f_2 \right| \mathrm{d}H^{\beta}_{\infty} > 0$, we must have,
			$$\frac{1}{H^{\beta}_{\infty} \left( Q' \right)} \int\limits_{Q'} \left| f_2 \right| \mathrm{d}H^{\beta}_{\infty} \leq \frac{K}{H^{\beta}_{\infty} \left( K \cdot Q \right)} \int\limits_{K \cdot Q} \left| f \right| \mathrm{d}H^{\beta}_{\infty} \leq K M_{\beta}f \left( x \right).$$
			Therefore,
			$$M_{\beta}f_2 \left( y \right) \leq K \ M_{\beta}f \left( x \right),$$
			for every $y \in Q$. This leads us to,
			\begin{equation}
				\label{F2Equation}
				\frac{1}{H^{\beta}_{\infty} \left( Q \right)} \int\limits_{Q} \left( M_{\beta}f_2 \right)^{\delta} \left( y \right) \mathrm{d}H^{\beta}_{\infty} \left( y \right) \leq K \left( M_{\beta}f \right)^{\delta} \left( x \right).
			\end{equation}
			Combining Equations \eqref{F1Equation} and \eqref{F2Equation} and using the fact that the Hausdorff content and the dyadic Hausdorff content are equivalent with the Choquet integral of the latter being sublinear, we get Equation \eqref{GoalEquation}. This completes the proof!
		\end{proof}

\end{document}
